% Setting up the document and importing necessary packages
\documentclass[12pt,a4paper]{article}

% Chinese support packages
\usepackage{xeCJK}
\usepackage{fontspec}
\setCJKmainfont{Noto Serif CJK TC}
\setCJKsansfont{Noto Sans CJK TC}
\setCJKmonofont{Noto Sans CJK TC}

% Other necessary packages
\usepackage{geometry}
\usepackage{titlesec}
\usepackage{enumitem}
\usepackage{color}
\usepackage{colortbl}
\usepackage{tcolorbox}
\usepackage{booktabs}
\usepackage{longtable}
\usepackage{array}
\usepackage{graphicx}
\usepackage{tikz}
\usepackage{mdframed}
\usepackage{fancyhdr}
\usepackage{hyperref}
\usepackage{multicol}
\usepackage{datetime2}
\usepackage{natbib} % For bibliography

\hypersetup{
    colorlinks=true,
    linkcolor=emergency_blue,
    citecolor=emergency_blue,
    urlcolor=emergency_blue,
    pdfborder={0 0 0}
}

% Page setup
\geometry{
    top=2.5cm,
    bottom=2.5cm,
    left=2cm,
    right=2cm,
    headheight=15pt
}

% Color definitions
\definecolor{emergency_red}{RGB}{186,24,27}
\definecolor{emergency_blue}{RGB}{0,114,188}
\definecolor{emergency_gray}{RGB}{120,120,120}
\definecolor{highlight_yellow}{RGB}{255,250,205}

% Title format settings
\titleformat{\section}[block]{\Large\bfseries\color{emergency_red}}{\thesection}{10pt}{}
\titleformat{\subsection}[block]{\large\bfseries\color{emergency_blue}}{\thesubsection}{8pt}{}
\titleformat{\subsubsection}[block]{\normalsize\bfseries\color{emergency_gray}}{\thesubsubsection}{6pt}{}

% Custom environment
\newmdenv[
    linecolor=emergency_red,
    backgroundcolor=highlight_yellow,
    linewidth=2pt,
    topline=true,
    bottomline=true,
    leftline=true,
    rightline=true,
    innertopmargin=10pt,
    innerbottommargin=10pt,
    innerleftmargin=10pt,
    innerrightmargin=10pt
]{importantbox}

\newcommand{\note}[1]{
    \par
    \noindent
    \begin{tcolorbox}[
        colback=emergency_blue!5!white,
        colframe=emergency_blue,
        title=\textbf{重點提醒},
        fonttitle=\bfseries,
        boxrule=0.5pt,
        arc=3pt,
        left=6pt,
        right=6pt,
        top=6pt,
        bottom=6pt
    ]
    #1
    \end{tcolorbox}
}

% Header and footer settings
\pagestyle{fancy}
\fancyhf{}
\fancyhead[L]{\textcolor{emergency_red}{內科部住院醫師教學文件}}
\fancyhead[R]{\textcolor{emergency_blue}{胸腔內科EPAs}}
\fancyfoot[C]{\textcolor{emergency_gray}{\thepage}}
\renewcommand{\headrulewidth}{1pt}
\renewcommand{\headrule}{\hbox to\headwidth{\color{emergency_red}\leaders\hrule height \headrulewidth\hfill}}

% Date format
\DTMsetstyle{yyyy-mm-dd}
\newcommand{\mydate}{\DTMtoday}

% Document start
\begin{document}

% Title page with table of contents
\begin{titlepage}
    \centering
    {\LARGE\textcolor{emergency_red}{\textbf{內科部住院醫師教學文件}}\par}
    \vspace{1cm}
    {\huge\bfseries 胸腔內科可信賴專業活動EPAs\par}
    \vspace{0.5cm}
    {\Large\textcolor{emergency_blue}{\textit{Pulmonology Entrustable Professional Activities}}\par}
    \vspace{1cm}
    
    \begin{tcolorbox}[
        colback=emergency_blue!5!white,
        colframe=emergency_blue,
        width=0.8\textwidth,
        arc=10pt,
        boxrule=1pt
    ]
    \centering
    {\large\textbf{目標對象}\par}
    \vspace{0.3cm}
    內科部住院醫師R1-R3\par
    胸腔內科專科訓練醫師\par
    \vspace{0.5cm}
    {\large\textbf{文件版本}\par}
    \vspace{0.3cm}
    第1.0版\par
    發布日期：\mydate\par
    \end{tcolorbox}

    \vfill
    
    {\large 溫岳峰 MD\\陳詩宇 MD \\謝慕揚 MD, PhD, FESC\\
    適用於CCC評量系統\par}
    
    \vspace{0.5cm}
\end{titlepage}

\newpage
\tableofcontents
\newpage

% Main content
\section{胸腔內科EPAs概述}

\subsection{EPA定義}
Entrustable Professional Activities (EPA) 是指可信賴委託給住院醫師執行的專業活動單元。胸腔內科EPAs涵蓋了呼吸系統疾病診斷、治療與照護的核心能力，是所有胸腔內科醫師必須精熟的專業技能。

\vspace{0.5cm}
\note{
胸腔內科EPAs評量著重於能否安全、穩定地執行完整的呼吸系統疾病診療流程，並能整合資訊做出適當臨床決策。每個EPA都包含明確的里程碑發展階段，從初學者到專家的完整能力建構。
}

\subsection{學習目標}
完成胸腔內科EPAs後，住院醫師應能夠：
\begin{itemize}[nosep]
    \item 熟練進行呼吸系統病史詢問與風險評估。
    \item 準確執行胸部理學檢查。
    \item 獨立判讀胸部X光與電腦斷層檢查。
    \item 診斷與處理急性呼吸衰竭。
    \item 管理各種氣道疾病。
    \item 治療肺炎與感染性疾病。
    \item 診斷與治療肺栓塞。
    \item 執行胸腔穿刺術與胸管置放。
    \item 操作支氣管內視鏡檢查。
    \item 協調胸腔科病患整體照護。
\end{itemize}

\subsection{委託等級}
\begin{table}[h]
    \centering
    \caption{EPA委託等級定義}
    \begin{tabular}{p{2cm} p{3cm} p{9cm}}
        \toprule
        \textbf{等級} & \textbf{委託程度} & \textbf{描述} \\
        \midrule
        Level 1 & 觀察學習 & 僅能觀察他人執行，無法穩立操作 \\
        Level 2 & 直接監督 & 可在主治醫師直接監督下執行 \\
        Level 3 & 間接監督 & 可穩立執行，但需回報並接受指導 \\
        Level 4 & 監督可得 & 可穩立執行，監督者隨時可得 \\
        Level 5 & 完全穩立 & 可穩立執行並指導他人 \\
        \bottomrule
    \end{tabular}
\end{table}

\newpage
\section{胸腔內科10大EPAs詳述}

\begin{longtable}{|p{1.5cm}|p{4cm}|p{8.5cm}|}
\caption{胸腔內科EPAs完整架構} \\
\hline
\rowcolor{emergency_blue!20}
\textbf{EPA} & \textbf{專業活動名稱} & \textbf{核心能力要素與里程碑} \\
\hline
\endfirsthead

\multicolumn{3}{c}%
{{\bfseries \tablename\ \thetable{} -- 接續上頁}} \\
\hline
\rowcolor{emergency_blue!20}
\textbf{EPA} & \textbf{專業活動名稱} & \textbf{核心能力要素與里程碑} \\
\hline
\endhead

\hline \multicolumn{3}{|r|}{{接續下頁}} \\ \hline
\endfoot

\hline \hline
\endlastfoot

\rowcolor{emergency_red!10}
\multicolumn{3}{|c|}{\textbf{EPA 1: 胸腔科病史詢問與風險評估}} \\
\hline
\textbf{描述} & \multicolumn{2}{|p{12.5cm}|}{能夠獨立進行呼吸系統疾病相關的完整病史詢問，包括症狀評估、危險因子識別、職業史調查，並能初步評估呼吸系統風險。} \\
\hline
\textbf{核心能力} & 症狀評估 & 咳嗽、咳痰、呼吸困難、胸痛、喘鳴的詳細詢問 \\
\cline{2-3}
& 危險因子識別 & 吸菸史、職業暴露、家族史、環境因子 \\
\cline{2-3}
& 功能狀態評估 & mMRC分級、運動耐受度評估 \\
\cline{2-3}
& 藥物史 & 呼吸系統藥物使用史、過敏史、順從性評估 \\
\cline{2-3}
& 職業與環境史 & 石棉、矽塵、有機粉塵暴露史 \\
\hline
\textbf{Level 1} & Novice & 能詢問基本呼吸系統症狀，需要監導者引導完成詢問 \\
\hline
\textbf{Level 2} & Advanced Beginner & 能系統性完成呼吸系統病史詢問，使用適當的症狀評估工具 \\
\hline
\textbf{Level 3} & Competent & 能獨立完成複雜病史詢問，整合症狀與危險因子進行初步風險分層 \\
\hline
\textbf{Level 4} & Proficient & 能在困難情況下獲取準確病史，整合多重共病資訊 \\
\hline
\textbf{Level 5} & Expert & 能指導他人進行病史詢問，識別罕見症狀表現 \\
\hline

\rowcolor{emergency_red!10}
\multicolumn{3}{|c|}{\textbf{EPA 2: 胸部理學檢查}} \\
\hline
\textbf{描述} & \multicolumn{2}{|p{12.5cm}|}{能夠熟練且準確地進行胸部的完整理學檢查，包括胸部視診、觸診、叩診、聽診，並能正確解讀檢查發現。} \\
\hline
\textbf{核心能力} & 胸部視診 & 胸廓形狀、呼吸型態、皮膚變化觀察 \\
\cline{2-3}
& 胸部觸診 & 胸廓擴張、語言震顫、皮下氣腫檢查 \\
\cline{2-3}
& 胸部叩診 & 肺界變化、濁音區域識別 \\
\cline{2-3}
& 胸部聽診 & 正常與異常呼吸音、囉音分類 \\
\cline{2-3}
& 其他檢查 & 杵狀指、發紺、頸靜脈怒張 \\
\hline
\textbf{Level 1} & Novice & 能進行基本胸部聽診，識別正常呼吸音 \\
\hline
\textbf{Level 2} & Advanced Beginner & 能識別常見異常呼吸音，進行系統性胸部檢查 \\
\hline
\textbf{Level 3} & Competent & 能區分各種異常囉音，準確評估胸腔積液程度，整合多項檢查發現 \\
\hline
\textbf{Level 4} & Proficient & 能識別複雜呼吸音特徵，檢查技術精準且有效率 \\
\hline
\textbf{Level 5} & Expert & 能偵測細微的異常發現，整合檢查發現與臨床表現 \\
\hline

\rowcolor{emergency_red!10}
\multicolumn{3}{|c|}{\textbf{EPA 3: 胸部影像判讀與分析}} \\
\hline
\textbf{描述} & \multicolumn{2}{|p{12.5cm}|}{能夠獨立判讀胸部X光與胸部電腦斷層，識別正常與異常變化，包括肺部病變、胸膜疾病、縱膈腔異常等，並提出臨床相關性分析。} \\
\hline
\textbf{核心能力} & 胸部X光判讀 & 肺野、心影、橫膈、骨骼結構評估 \\
\cline{2-3}
& 胸部CT判讀 & 肺實質、支氣管、血管、淋巴結評估 \\
\cline{2-3}
& 病變識別 & 結節、腫塊、浸潤、纖維化變化 \\
\cline{2-3}
& 胸膜疾病 & 胸膜積液、氣胸、胸膜增厚評估 \\
\cline{2-3}
& 縱膈腔評估 & 淋巴結腫大、縱膈腔腫瘤識別 \\
\hline
\textbf{Level 1} & Novice & 能識別基本正常胸部X光結構 \\
\hline
\textbf{Level 2} & Advanced Beginner & 能識別明顯異常病變，判讀基本胸部CT \\
\hline
\textbf{Level 3} & Competent & 能獨立判讀大部分胸部影像異常，整合影像與臨床症狀 \\
\hline
\textbf{Level 4} & Proficient & 能判讀複雜影像變化，識別細微的早期病變 \\
\hline
\textbf{Level 5} & Expert & 成為影像判讀的專家，能教導複雜影像判讀 \\
\hline

\rowcolor{emergency_red!10}
\multicolumn{3}{|c|}{\textbf{EPA 4: 肺功能檢查操作與判讀}} \\
\hline
\textbf{描述} & \multicolumn{2}{|p{12.5cm}|}{能夠進行基礎肺功能檢查，包括肺量計操作、品質控制、基本測量與主要結果評估，並能識別重要的功能異常模式。} \\
\hline
\textbf{核心能力} & 檢查操作 & 肺量計操作、病患指導、品質確保 \\
\cline{2-3}
& 基本測量 & FVC、FEV1、FEV1/FVC比值測量 \\
\cline{2-3}
& 功能評估 & 阻塞性、限制性功能異常識別 \\
\cline{2-3}
& 支氣管擴張試驗 & 試驗執行與結果判讀 \\
\cline{2-3}
& 臨床關聯 & 功能異常與疾病嚴重度關聯 \\
\hline
\textbf{Level 1} & Novice & 能協助進行基本肺量計檢查 \\
\hline
\textbf{Level 2} & Advanced Beginner & 能獨立操作標準肺功能檢查，進行基本判讀 \\
\hline
\textbf{Level 3} & Competent & 能完成標準肺功能檢查，準確評估功能異常類型 \\
\hline
\textbf{Level 4} & Proficient & 能進行進階肺功能技術，準確定量評估 \\
\hline
\textbf{Level 5} & Expert & 成為肺功能技術專家，能教導複雜技術 \\
\hline

\rowcolor{emergency_red!10}
\multicolumn{3}{|c|}{\textbf{EPA 5: 急性呼吸衰竭診斷與初步處理}} \\
\hline
\textbf{描述} & \multicolumn{2}{|p{12.5cm}|}{能夠快速識別急性呼吸衰竭，進行適當的診斷性檢查，評估嚴重度分級，並執行初步治療措施，包括氧氣治療與呼吸器支持決策。} \\
\hline
\textbf{核心能力} & 快速識別 & 急性呼吸窘迫症狀識別、高風險特徵 \\
\cline{2-3}
& 診斷策略 & 動脈血氣體、胸部影像、實驗室檢查安排 \\
\cline{2-3}
& 嚴重度分級 & P/F ratio、SOFA評分應用 \\
\cline{2-3}
& 氧氣治療 & 鼻導管、面罩、高流量氧氣選擇 \\
\cline{2-3}
& 呼吸器決策 & NIPPV適應症、插管時機、轉院必要性 \\
\hline
\textbf{Level 1} & Novice & 能識別典型呼吸窘迫症狀，知道基本檢查項目 \\
\hline
\textbf{Level 2} & Advanced Beginner & 能使用標準診斷流程，進行基本嚴重度評估 \\
\hline
\textbf{Level 3} & Competent & 能獨立診斷急性呼吸衰竭，適當選擇治療策略 \\
\hline
\textbf{Level 4} & Proficient & 能處理複雜臨床情況，整合多重因子做決策 \\
\hline
\textbf{Level 5} & Expert & 成為急性呼吸衰竭專家，能指導複雜個案處理 \\
\hline

\rowcolor{emergency_red!10}
\multicolumn{3}{|c|}{\textbf{EPA 6: 氣道疾病診斷與處理}} \\
\hline
\textbf{描述} & \multicolumn{2}{|p{12.5cm}|}{能夠診斷各種氣道疾病，評估氣流阻塞程度，選擇適當的治療方式，包括支氣管擴張劑、類固醇治療，或轉介進階治療。} \\
\hline
\textbf{核心能力} & 疾病識別 & 氣喘、COPD、支氣管擴張症分類 \\
\cline{2-3}
& 嚴重度評估 & 症狀控制、急性惡化評估 \\
\cline{2-3}
& 治療選擇 & 吸入型藥物vs全身性治療決策 \\
\cline{2-3}
& 急性處置 & 急性氣喘發作、COPD急性惡化處理 \\
\cline{2-3}
& 長期管理 & 控制型藥物、肺復健適應症 \\
\hline
\textbf{Level 1} & Novice & 能識別基本氣道阻塞症狀，知道緊急處置原則 \\
\hline
\textbf{Level 2} & Advanced Beginner & 能處理穩定性氣道疾病，選擇基本吸入型藥物 \\
\hline
\textbf{Level 3} & Competent & 能獨立處理大部分氣道疾病，做出急性惡化決定 \\
\hline
\textbf{Level 4} & Proficient & 能處理複雜氣道疾病，評估生物製劑適應症 \\
\hline
\textbf{Level 5} & Expert & 成為氣道疾病專家，能執行複雜治療程序 \\
\hline

\rowcolor{emergency_red!10}
\multicolumn{3}{|c|}{\textbf{EPA 7: 肺炎與感染性疾病診斷與治療}} \\
\hline
\textbf{描述} & \multicolumn{2}{|p{12.5cm}|}{能夠診斷各種肺炎與呼吸道感染，區分社區型與院內感染，制定個別化治療計畫，包括抗生素選擇與非藥物治療。} \\
\hline
\textbf{核心能力} & 診斷評估 & 症狀、徵象、嚴重度分級評估 \\
\cline{2-3}
& 病因識別 & 典型vs非典型病原菌區分 \\
\cline{2-3}
& 嚴重度評估 & CURB-65、PSI評分應用 \\
\cline{2-3}
& 抗生素治療 & 經驗性與標靶治療選擇 \\
\cline{2-3}
& 整合照護 & 住院指標、疫苗接種建議 \\
\hline
\textbf{Level 1} & Novice & 能識別肺炎症狀徵象，知道基本抗生素治療 \\
\hline
\textbf{Level 2} & Advanced Beginner & 能使用標準治療指引，進行基本抗生素調整 \\
\hline
\textbf{Level 3} & Competent & 能獨立管理穩定肺炎，調整複雜抗生素組合 \\
\hline
\textbf{Level 4} & Proficient & 能處理難治性感染，評估進階治療選項 \\
\hline
\textbf{Level 5} & Expert & 成為感染疾病專家，發展個人化治療策略 \\
\hline

\rowcolor{emergency_red!10}
\multicolumn{3}{|c|}{\textbf{EPA 8: 肺栓塞診斷與管理}} \\
\hline
\textbf{描述} & \multicolumn{2}{|p{12.5cm}|}{能夠診斷肺栓塞，評估血栓栓塞風險，制定個別化抗凝血治療策略，監測治療效果與副作用，並提供預防性建議。} \\
\hline
\textbf{核心能力} & 臨床評估 & Wells評分、Geneva評分應用 \\
\cline{2-3}
& 診斷策略 & D-dimer、CTPA、V/Q scan選擇 \\
\cline{2-3}
& 風險分層 & 血流動力學影響評估 \\
\cline{2-3}
& 抗凝治療 & DOAC vs warfarin選擇 \\
\cline{2-3}
& 預防措施 & VTE預防、長期管理 \\
\hline
\textbf{Level 1} & Novice & 能識別肺栓塞症狀，知道基本診斷流程 \\
\hline
\textbf{Level 2} & Advanced Beginner & 能使用臨床評分工具，開始抗凝血治療 \\
\hline
\textbf{Level 3} & Competent & 能獨立診斷肺栓塞，選擇適當抗凝治療 \\
\hline
\textbf{Level 4} & Proficient & 能處理複雜肺栓塞，評估血栓溶解治療 \\
\hline
\textbf{Level 5} & Expert & 成為血栓栓塞專家，發展個人化治療策略 \\
\hline

\rowcolor{emergency_red!10}
\multicolumn{3}{|c|}{\textbf{EPA 9: 胸腔穿刺術與胸管置放}} \\
\hline
\textbf{描述} & \multicolumn{2}{|p{12.5cm}|}{能夠執行胸腔穿刺術與胸管置放術，包括適應症評估、解剖定位、無菌技術、併發症預防，並能處理術後照護與管理。} \\
\hline
\textbf{核心能力} & 適應症評估 & 診斷性vs治療性穿刺決策 \\
\cline{2-3}
& 解剖定位 & 超音波導引、解剖標誌識別 \\
\cline{2-3}
& 無菌技術 & 消毒、局部麻醉、穿刺技巧 \\
\cline{2-3}
& 併發症預防 & 氣胸、出血、感染預防 \\
\cline{2-3}
& 術後管理 & 胸管照護、移除時機 \\
\hline
\textbf{Level 1} & Novice & 能協助胸腔穿刺，了解基本解剖結構 \\
\hline
\textbf{Level 2} & Advanced Beginner & 能在監督下執行診斷性胸腔穿刺 \\
\hline
\textbf{Level 3} & Competent & 能獨立執行胸腔穿刺，處理一般併發症 \\
\hline
\textbf{Level 4} & Proficient & 能執行複雜胸腔手術，使用進階技術 \\
\hline
\textbf{Level 5} & Expert & 成為胸腔手術專家，能教導複雜技術 \\
\hline

\rowcolor{emergency_red!10}
\multicolumn{3}{|c|}{\textbf{EPA 10: 支氣管內視鏡檢查}} \\
\hline
\textbf{描述} & \multicolumn{2}{|p{12.5cm}|}{能夠執行基礎支氣管內視鏡檢查，包括設備操作、解剖識別、檢體取得、治療性處置，並能識別重要的病理發現。} \\
\hline
\textbf{核心能力} & 設備操作 & 支氣管鏡操作、影像調整 \\
\cline{2-3}
& 解剖識別 & 氣管支氣管樹解剖認知 \\
\cline{2-3}
& 檢體取得 & 沖洗液、刷檢、切片技術 \\
\cline{2-3}
& 病理識別 & 發炎、腫瘤、異物識別 \\
\cline{2-3}
& 治療處置 & 止血、異物移除、支架置放 \\
\hline
\textbf{Level 1} & Novice & 能協助支氣管鏡檢查，識別基本解剖 \\
\hline
\textbf{Level 2} & Advanced Beginner & 能在監督下執行診斷性支氣管鏡 \\
\hline
\textbf{Level 3} & Competent & 能獨立執行標準支氣管鏡，取得適當檢體 \\
\hline
\textbf{Level 4} & Proficient & 能執行治療性支氣管鏡，處理併發症 \\
\hline
\textbf{Level 5} & Expert & 成為支氣管鏡專家，執行進階技術 \\
\hline

\end{longtable}

\newpage
\section{評量方法與標準}

\subsection{CCC評量架構}
本胸腔內科EPAs採用Case-based Clinical Competency (CCC)評量方式：
\begin{itemize}[nosep]
    \item 真實臨床情境評量
    \item 多次觀察與回饋
    \item 漸進式委託決策
    \item 360度回饋機制
\end{itemize}

\begin{importantbox}
\textbf{特別提醒：} 胸腔內科EPAs的評量重點在於漸進式能力發展，而非單次完美表現。鼓勵從錯誤中學習，持續改進。每個EPA都需要在真實臨床情境中反復練習與評量。
\end{importantbox}

\subsection{評量工具對應表}
\begin{table}[h]
    \centering
    \caption{各EPA推薦評量工具}
    \begin{tabular}{p{4cm} p{4cm} p{6cm}}
        \toprule
        \textbf{EPA類型} & \textbf{主要評量工具} & \textbf{輔助評量方法} \\
        \midrule
        技能操作類 & DOPS & 直接觀察、技能檢核表 \\
        (EPA 4, 9, 10) & (Direct Observation of Procedural Skills) & 同儕評量、自我評量 \\
        \midrule
        臨床推理類 & Mini-CEX & 病例討論、口試 \\
        (EPA 1, 3, 5, 6, 7, 8) & (Mini-Clinical Evaluation Exercise) & 病歷撰寫評量、模擬情境 \\
        \midrule
        理學檢查類 & DOPS & 標準化病人、技能測驗 \\
        (EPA 2) & (Direct Observation of Procedural Skills) & 影片回顧、同儕觀察 \\
        \bottomrule
    \end{tabular}
\end{table}

\subsection{評量頻率建議}
\begin{itemize}[nosep]
    \item \textbf{每月評量}：選擇2-3個EPA重點評量
    \item \textbf{季度檢討}：全面檢視所有EPA進展
    \item \textbf{年度認證}：EPA達標認證與進階規劃
    \item \textbf{即時回饋}：發現學習需求時立即介入
\end{itemize}

\vspace{0.5cm}
\note{
胸腔內科EPAs的學習是一個持續的過程。建議住院醫師建立個人學習檔案，記錄每次評量結果與改善計畫。主治醫師應提供及時具體的回饋，協助學習者達到下一個里程碑。
}

\section{實施建議與學習資源}

\subsection{月度晨會Mini-OSCE整合建議}
\begin{table}[h]
    \centering
    \caption{月度EPA評量主題安排}
    \begin{tabular}{p{3cm} p{5cm} p{6cm}}
        \toprule
        \textbf{月份} & \textbf{重點EPAs} & \textbf{評量重點} \\
        \midrule
        第1個月 & EPA 1, 2, 3 & 基礎評估技能建立 \\
        第2個月 & EPA 4, 5, 6 & 診斷技術與急症處理 \\
        第3個月 & EPA 7, 8 & 感染與血管疾病管理 \\
        第4個月 & EPA 9, 10, 綜合評量 & 進階技能與整合能力 \\
        \bottomrule
    \end{tabular}
\end{table}

\subsection{推薦學習資源}
學員可參考以下文獻深入學習：
\begin{itemize}
    \item Murray \& Nadel's Textbook of Respiratory Medicine. 7th edition.
    \item GOLD Guidelines for COPD Management (2024 Update)
    \item ATS/ERS Guidelines for Pulmonary Function Testing
    \item BTS Guidelines for Management of Community Acquired Pneumonia
\end{itemize}

\subsection{自我評估工具}
\begin{itemize}[nosep]
    \item 定期記錄學習歷程與反思
    \item 尋求多方回饋（主治醫師、護理師、病人）
    \item 參與同儕學習與討論
    \item 利用模擬教學工具練習
    \item 建立個人學習檔案
\end{itemize}

\section{總結}

胸腔內科EPAs涵蓋了呼吸系統疾病診療的核心能力，從基礎的病史詢問與理學檢查，到複雜的支氣管內視鏡檢查與胸腔手術。每個EPA都有清楚的能力描述與里程碑發展路徑，協助住院醫師循序漸進地建立專業能力。

\begin{importantbox}
\textbf{關鍵訊息：} 胸腔內科EPAs的核心在於「可委託性」- 當您能夠安全、穩立地執行各項呼吸系統疾病診療活動，並做出合理的臨床決策時，就達到了相對應EPA的目標。記住，每一次的臨床實務都是學習的機會，透過持續的實作、反思與改進，才能成為一位優秀的胸腔內科醫師。
\end{importantbox}

% References section
\section{參考文獻}
\begin{thebibliography}{9}

\bibitem{murray2022}
Broaddus, V. C., Ernst, J. D., King Jr, T. E., Lazarus, S. C., Gotway, M. B., Sarmiento, K. F., \& Wolters, P. J. (2022). \textit{Murray \& Nadel's Textbook of Respiratory Medicine}. 7th edition. Elsevier.

\bibitem{gold2024}
Global Initiative for Chronic Obstructive Lung Disease. (2024). Global Strategy for the Diagnosis, Management, and Prevention of Chronic Obstructive Pulmonary Disease. \textit{GOLD Report}.

\bibitem{olle2013}
ten Cate, O. (2013). Nuts and bolts of entrustable professional activities. \textit{Journal of Graduate Medical Education}, 5(1), 157-158.

\bibitem{ats2019}
Graham, B. L., et al. (2019). Standardization of Spirometry 2019 Update. \textit{American Journal of Respiratory and Critical Care Medicine}, 200(8), e70-e88.

\bibitem{bts2022}
Lim, W. S., et al. (2022). BTS guidelines for the management of community acquired pneumonia in adults: update 2022. \textit{Thorax}, 77(8), 717-724.

\end{thebibliography}

\end{document}